\ifdefined\gkver\else\def\gkver{student}\fi
\documentclass[\gkver]{gaokaozhenti}
\begin{document}

\chapter{2000年全国旧课程}

\begin{enumerate}
\item
%% number: 1
%% paperName: 2000年全国高考旧课程第 1 题 4 分
%% typeId: 1
%% type: 单选题
%% chapter: 原子和原子核
%% point: 天然放射性
%% method: 无
%% score: 4
%% degree: 850
%% duplicateId: 0
%% body:
最近几年，原子核科学家在超重元素岛的探测方面取得重大进展。1996 年科学家们在研究某两个重离子结合成超重元素的反应时，发现生成的超重元素的核 $\ce{^{A}_{Z}X}$ 经过 6 次 $\alpha$ 衰变后的产物是 $\ce{^{253}_{100}Fm}$。由此，可判定生成的超重元素的原子序数和质量数分别是 \xzanswer{D}

\fourchoices[answer=D]
{124，259}
{124，265}
{112，265}
{112，277}

%% answer: D
%% memo:
\memoanswer{
核反应方程满足质量数和核电荷数守恒。每经过一次 $\alpha$ 衰变质量数减少 $4$，有 $A = 253 + 4 \times 6 = 277$；又每经过一次 $\alpha$ 衰变核电荷数减少 $2$，每经过一次 $\beta$ 衰变核电荷数增加 $1$，有 $Z = 100 + 6 \times 2 = 112$。超重元素的原子序数为 $112$，质量数为 $277$，故 D 正确。
}

\item
%% number: 2
%% paperName: 2000年全国高考旧课程第 2 题 4 分
%% typeId: 1
%% type: 单选题
%% chapter: 气体性质
%% point: 理想气体状态方程
%% method: 无
%% score: 4
%% degree: 850
%% duplicateId: 0
%% body:
对于一定量的理想气体，下列四个论述中正确的是 \xzanswer{B}

\fourchoices[answer=B]
{当分子热运动变剧烈时，压强必变大}
{当分子热运动变剧烈时，压强可以不变}
{当分子间的平均距离变大时，压强必变小}
{当分子间的平均距离变大时，压强必变大}

%% answer: B
%% memo:
\memoanswer{
影响理想气体的压强的因素有温度、体积。由于热运动都加强，温度升高，但无法确定气体体积的变化，压强的变化就不一定，故 A 错误，B 正确；由 C、D 两项可判断气体体积变大，但温度的变化是不确定的，压强的变化就不是确定的，故 C、D 错误。
}

\item
%% number: 3
%% paperName: 2000年全国高考旧课程第 3 题 4 分
%% typeId: 1
%% type: 单选题
%% chapter: 曲线运动
%% point: 人造地球卫星
%% method: 无
%% score: 4
%% degree: 550
%% duplicateId: 0
%% body:
某人造地球卫星因受高空稀薄空气的阻力作用，绕地球运转的轨道会慢慢改变，每次测量中卫星的运动可近似看作圆周运动。某次测量卫星的轨道半径为 $r_{1}$，后来变为 $r_{2}$，$r_{2} < r_{1}$。以 $E_{k1}$、$E_{k2}$ 表示卫星在这两个轨道上的动能，$T_{1}$、$T_{2}$ 表示卫星在这两个轨道上绕地运动的周期，则 \xzanswer{C}

\fourchoices[answer=C]
{$E_{k2} < E_{k1}$，$T_{2} < T_{1}$}
{$E_{k2} < E_{k1}$，$T_{2} > T_{1}$}
{$E_{k2} > E_{k1}$，$T_{2} < T_{1}$}
{$E_{k2} > E_{k1}$，$T_{2} > T_{1}$}

%% answer: C
%% memo:
\memoanswer{
人造卫星在某轨道环绕时，其万有引力提供向心力，有 $G\frac{Mm}{r^{2}}=m\cdot\frac{4\pi^{2}}{T^{2}}r$ 或 $G\frac{Mm}{r^{2}}=\frac{mv^{2}}{r}$，从上面两式可以看出 $T^{2}\propto r^{3}$ 或 $\frac{1}{2}mv^{2}\propto\frac{1}{r}$，故 C 正确。
}

\item
%% number: 4
%% paperName: 2000年全国高考旧课程第 4 题 4 分
%% typeId: 2
%% type: 多选题
%% chapter: 电路
%% point: 电容
%% method: 无
%% score: 4
%% degree: 650
%% duplicateId: 0
%% body:
对于水平放置的平行板电容器，下列说法正确的是 \xzanswer{BCD}

\fourchoices[answer=BCD]
{将两极板的间距加大，电容将增大}
{将两极板平行错开，使正对面积减小，电容将减小}
{在下板的内表面上放置一面积和极板相等、厚度小于极板间距的陶瓷板，电容将增大}
{在下板的表面上放置一面积和极板相等、厚度小于极板间距的铝板，电容将增大}

%% answer: BCD
%% memo:
\memoanswer{
对于平行板电容器有 $C=\frac{\varepsilon S}{4\pi kd}$，可知 A 错误，B、C 正确；当插入金属板时，由静电感应，相当于电容器的两板距离变小，故 D 正确。
}

\item
%% number: 5
%% paperName: 2000年全国高考旧课程第 5 题 4 分
%% typeId: 1
%% type: 单选题
%% chapter: 气体性质
%% point: 热力学第一定律
%% method: 无
%% score: 4
%% degree: 550
%% duplicateId: 0
%% body:
图中活塞将气缸分成甲、乙两气室，气缸、活塞（连同拉杆）是绝热的，且不漏气。以 $E_{\text{甲}}$、$E_{\text{乙}}$ 分别表示甲、乙两气室中气体的内能，则在将拉杆缓慢向外拉的过程中 \xzanswer{C}
\onepicture[w=6cm]{figs/05.png}

\fourchoices[answer=C]
{$E_{\text{甲}}$ 不变，$E_{\text{乙}}$ 减小}
{$E_{\text{甲}}$ 增大，$E_{\text{乙}}$ 不变}
{$E_{\text{甲}}$ 增大，$E_{\text{乙}}$ 减小}
{$E_{\text{甲}}$ 不变，$E_{\text{乙}}$ 不变}

%% answer: C
%% memo:
\memoanswer{
因气缸绝热，当活塞缓慢向外拉时，甲、乙两团气体与外界不发生热量交换，对甲气体活塞对它做功，内能将增大，而对乙气体则对外做功，内能将减小，故 C 正确。
}

\item
%% number: 6
%% paperName: 2000年全国高考旧课程第 6 题 4 分
%% typeId: 1
%% type: 单选题
%% chapter: 光学
%% point: 光的电磁说
%% method: 无
%% score: 4
%% degree: 650
%% duplicateId: 0
%% body:
图为 X 射线管的结构示意图，E 为灯丝电源，要使射线管发出 X 射线，须在 K、A 两电极间加上几万伏的直流高压 \xzanswer{D}
\onepicture[h=4.2cm]{figs/06.png}

\fourchoices[answer=D]
{高压电源正极应接在 P 点，X 射线从 K 极发出}
{高压电源正极应接在 P 点，X 射线从 A 极发出}
{高压电源正极应接在 Q 点，X 射线从 K 极发出}
{高压电源正极应接在 Q 点，X 射线从 A 极发出}

%% answer: D
%% memo:
\memoanswer{
伦琴射线是高速电子流撞到固体的表面而形成的。在原题图中，$K$ 为阴极，其作用是发射电子，$A$ 是阳极，也叫对阴极，在 $A$、$K$ 间加上高电压，将电子从 $K$ 加速到 $A$，高速电子撞到 $A$ 后，辐射出 X 射线，故 D 正确。
}

\item
%% number: 7
%% paperName: 2000年全国高考旧课程第 7 题 4 分
%% typeId: 2
%% type: 多选题
%% chapter: 机械波
%% point: 波的多解性
%% method: 无
%% score: 4
%% degree: 650
%% duplicateId: 0
%% body:
一列横波在 $t=0$ 时刻的波形如图中实线所示，在 $t=1\Us$ 时刻的波形如图中虚线所示，由此可以判定此波的 \xzanswer{AC}
\onepicture[w=7.5cm]{figs/07.png}

\fourchoices[answer=AC]
{波长一定是 $4\Ucm$}
{周期一定是 $4\Us$}
{振幅一定是 $2\Ucm$}
{传播速度一定是 $1\Ucms$}

%% answer: AC
%% memo:
\memoanswer{
从图中可直接确定波长为 $4\Ucm$ 和振幅 $2\Ucm$，故 A、C 正确；由于不知波的传播方向，所以周期和波速不能确定。由图可知 $\lambda = 4\Ucm$，$A = 2\Ucm$，但其周期和波速有多组解。

若横波向 $x$ 正方向传播，则 $\left(n+\frac{1}{4}\right)T=1\Us$，解得 $T=\frac{4}{4n+1}\Us$（$n=0,1,2,3,\dots$），有 $v=\frac{\lambda}{T}=(4n+1)\Ucms$（$n=0,1,2,3,\dots$）；若横波向 $x$ 负方向传播，则 $\left(n+\frac{3}{4}\right)T=1\Us$，解得 $T=\frac{4}{4n+3}\Us$（$n=0,1,2,3,\dots$），有 $v=\frac{\lambda}{T}=(4n+3)\Ucms$（$n=0,1,2,3,\dots$），故 B、D 错误。
}

\item
%% number: 8
%% paperName: 2000年全国高考旧课程第 8 题 4 分
%% typeId: 2
%% type: 多选题
%% chapter: 光学
%% point: 光的折射
%% method: 无
%% score: 4
%% degree: 650
%% duplicateId: 0
%% body:
A 与 B 是两束平行的单色光，它们从空气射入水中的折射角分别为 $r_{A}$、$r_{B}$，若 $r_{A} > r_{B}$，则 \xzanswer{AB}

\fourchoices[answer=AB]
{在空气中 A 波长大于 B 的波长}
{在水中 A 的传播速度大于 B 的传播速度}
{A 的频率大于 B 的频率}
{在水中 A 的波长小于 B 的波长}

%% answer: AB
%% memo:
\memoanswer{
因 $r_{A} > r_{B}$，所以 $n_{A} < n_{B}$，$f_{A} < f_{B}$。在空气中光速相同，由 $\lambda=\frac{c}{f}$ 知 $\lambda_{A}>\lambda_{B}$，故 A 正确；又因为 $v=\frac{c}{n}$，在水中 A 光的传播速度大于 B 的传播速度，故 B 正确，C、D 错误。
}

\item
%% number: 9
%% paperName: 2000年全国高考旧课程第 9 题 4 分
%% typeId: 2
%% type: 多选题
%% chapter: 电路
%% point: 故障分析
%% method: 无
%% score: 4
%% degree: 650
%% duplicateId: 0
%% body:
图为一电路板的示意图，a、b、c、d 为接线柱，a、b 与 $200\UV$ 的交流电源连接，ab 间、bc 间、cd 间分别连接一个电阻。现发现电路中没有电流，为检查电路故障，有一交流电压表分别测得 b、d 两点间以及 a、c 两点间的电压均为 $200\UV$。由此可知 \xzanswer{CD}
\onepicture[w=4.5cm]{figs/09.png}

\fourchoices[answer=CD]
{ab 间电路通，cd 间电路不通}
{ab 间电路不通，bc 间电路通}
{ab 间电路通，bc 间电路不通}
{bc 间电路不通，cd 间电路通}

%% answer: CD
%% memo:
\memoanswer{
从原题的图中可以看出，该电路是一个串联电路。因为电路中无电流，所以故障一定是某处断路。根据两次测量，断路故障既处在 $bd$ 间，又处在 $ac$ 间，断路处在这两部分的公共部分，即 $bc$ 部分，故 C、D 正确。
}

\item
%% number: 10
%% paperName: 2000年全国高考旧课程第 10 题 4 分
%% typeId: 1
%% type: 单选题
%% chapter: 直线运动
%% point: 运动的合成
%% method: 无
%% score: 4
%% degree: 450
%% duplicateId: 0
%% body:
图为空间探测器的示意图，$P_{1}$、$P_{2}$、$P_{3}$、$P_{4}$ 是四个喷气发动机，$P_{1}$、$P_{3}$ 的连线与空间一固定坐标系的 $x$ 轴平行，$P_{2}$、$P_{4}$ 的连线与 $y$ 轴平行。每台发动机开动时，都能向探测器提供推力，但不会使探测器转动。开始时，探测器以恒定的速率 $v_{0}$ 向正 $x$ 方向平动，要使探测器改为向正 $x$ 偏负 $y$ $60^\circ$ 的方向以原来的速率 $v_{0}$ 平动，则可 \xzanswer{A}
\onepicture[w=5.5cm]{\includesvg{figs/10}}

\fourchoices[answer=A]
{先开动 $P_{1}$ 适当时间，再开动 $P_{4}$ 适当时间}
{先开动 $P_{3}$ 适当时间，再开动 $P_{2}$ 适当时间}
{开动 $P_{4}$ 适当时间}
{先开动 $P_{3}$ 适当时间，再开动 $P_{4}$ 适当时间}

%% answer: A
%% memo:
\memoanswer{
$P_{1}$ 发动时，使探测器得到 $-x$ 方向的速度；$P_{4}$ 发动时，使探测器得到 $-y$ 方向的速度。要使探测器沿 $+x$ 偏 $-y$ $60^\circ$ 方向以原来的速率 $v_{0}$ 平动，就必须使 $+x$ 方向速度减小，使 $-y$ 方向速度增加，所以要开动 $P_{1}$ 适当时间，再开动 $P_{4}$ 适当时间，故 A 正确。
}

\item
%% number: 11
%% paperName: 2000年全国高考旧课程第 11 题 5 分
%% typeId: 3
%% type: 填空题
%% chapter: 原子和原子核
%% point: 核裂变
%% method: 无
%% score: 5
%% degree: 650
%% duplicateId: 0
%% body:
裂变反应是目前利用中常用的反应，以原子核 $\ce{^{235}_{92}U}$ 为燃料的反应堆中，当 $\ce{^{235}_{92}U}$ 俘获一个慢中子后发生的裂变反应可以有多种方式，其中一种可表示为

$\ce{^{235}_{92}U + ^{1}_{0}n -> ^{139}_{54}Xe + ^{94}_{38}Sr + 3^{1}_{0}n}$

$235.0439 \qquad 1.0087 \qquad 138.9178 \qquad 93.9154$

反应方程下方的数字为中子及有关原子的静止质量（以原子质量单位 $\Uu$ 为单位）。已知 $1\Uu$ 的质量对应的能量为 $9.3\times10^{2}\UMeV$，此裂变反应释放出的能量是 \tkanswer[3cm]{$1.8\times10^{2}$} $\UMeV$。

%% answer:
\jdanswer{
$1.8\times10^{2}$
}

%% memo:
\memoanswer{
由题意可计算核反应前后的质量亏损为 $\Delta m = 235.0439 + 1.0087 - 138.9178 - 93.9154 - 3 \times 1.0087 = 0.1933\Uu$，故此裂变释放出的能量为 $\Delta E = \Delta m c^{2} = 0.1933 \times 9.3 \times 10^{2}\UMeV = 1.8 \times 10^{2}\UMeV$。

评分标准：本题 5 分，答案正确给 5 分，答错或不答给 0 分。
}

\item
%% number: 12
%% paperName: 2000年全国高考旧课程第 12 题 5 分
%% typeId: 3
%% type: 填空题
%% chapter: 电磁感应
%% point: 线圈过有界磁场
%% method: 无
%% score: 5
%% degree: 450
%% duplicateId: 0
%% body:
空间存在以 ab、cd 为边界的匀强磁场区域，磁感应强度大小为 $B$，方向垂直纸面向外，区域宽为 $l_{1}$。现有一矩形线框处在图中纸面内，它的短边与 ab 重合，长度为 $l_{2}$，长边的长度为 $2l_{1}$，如图所示。某时刻线框以初速度 $v$ 沿与 ab 垂直的方向进入磁场区域，同时某人对线框施以作用力，使它的速度大小和方向保持不变。设该线框的电阻为 $R$。从线框开始进入磁场到完全离开磁场的过程中，人对线框作用力所做的功等于 \tkanswer[4cm]{$\frac{2B^{2}l_{2}^{2}l_{1}v}{R}$}。
\onepicture[h=3.6cm, align=right]{figs/12.png}

%% answer:
\jdanswer{
$\frac{2B^{2}l_{2}^{2}l_{1}v}{R}$
}

%% memo:
\memoanswer{
由题意知，只有线框右边进入磁场和左边离开磁场的过程中，线框中才有感应电流，外力才做功；当左右两边都不在磁场中时，线框中无感应电流，外力不做功。由功能关系得 $W = 2BIl_{2}\cdot l_{1} = 2Bl_{1}l_{2}\cdot\frac{Bl_{2}v}{R} = \frac{2B^{2}l_{1}l_{2}^{2}v}{R}$。
}

\item
%% number: 13
%% paperName: 2000年全国高考旧课程第 13 题 5 分
%% typeId: 3
%% type: 填空题
%% chapter: 原子和原子核
%% point: 原子核式结构
%% method: 无
%% score: 5
%% degree: 450
%% duplicateId: 0
%% body:
假设在 $\ce{NaCl}$ 蒸气中存在由钠离子和氯离子靠静电相互作用构成的单个氯化钠 $\ce{NaCl}$ 分子。若取 $\ce{Na+}$ 与 $\ce{Cl-}$ 相距无限远时其电势能为零，一个 $\ce{NaCl}$ 分子的电势能为 $-6.1\UeV$。已知使一个中性钠原子 $\ce{Na}$ 最外层的电子脱离钠原子而形成钠离子所需的能量（电离能）为 $5.1\UeV$，使一个中性氯原子 $\ce{Cl}$ 结合一个电子形成氯离子所放出的能量（亲和能）为 $3.8\UeV$。由此可算出，在将一个 $\ce{NaCl}$ 分子分解成彼此远离的中性钠原子 $\ce{Na}$ 和中性氯原子 $\ce{Cl}$ 的过程中，外界供给的总能量等于 \tkanswer[2cm]{4.8} $\UeV$。

%% answer:
\jdanswer{
$4.8$
}

%% memo:
\memoanswer{
由题意可知，将 $\ce{NaCl}$ 变成 $\ce{Na+}$ 和 $\ce{Cl-}$ 吸收能量为 $6.1\UeV$，再将 $\ce{Na+}$ 变成 $\ce{Na}$ 放出 $5.1\UeV$ 的能量，$\ce{Cl-}$ 离子变成氯原子 $\ce{Cl}$ 吸收 $3.8\UeV$，所以将 $\ce{NaCl}$ 分解成中性的钠 $\ce{Na}$ 和中性的氯原子 $\ce{Cl}$ 的过程中，外界供给的能量等于 $6.1 + 3.8 - 5.1 = 4.8\UeV$。

正负电荷之间的相对位置的变化，引起两个电荷的电势能的变化。当距离变大时，电场力做负功，需要吸收能量：先把 $\ce{NaCl}$ 变成 $\ce{Na+}$ 和 $\ce{Cl-}$ 吸收能量为 $6.1\UeV$，再将 $\ce{Na+}$ 变成 $\ce{Na}$ 的过程放出能量 $5.1\UeV$，同时将 $\ce{Cl-}$ 变成 $\ce{Cl}$，电场力做负功电势能增大，也就是吸收能量。对于整个过程可以看成是由两个吸收能量和一个放出能量的过程组成的。

评分标准：本题 5 分，答案正确给 5 分，答错或不答给 0 分。
}

\item
%% number: 14
%% paperName: 2000年全国高考旧课程第 14 题 6 分
%% typeId: 4
%% type: 实验题
%% chapter: 运动和力
%% point: 验证动量守恒实验
%% method: 无
%% score: 6
%% degree: 450
%% duplicateId: 0
%% body:
某同学用图\ref{2000全国旧课程14a}所示装置通过半径相同的 A、B 两球的碰撞来验证动量守恒定律。图中 PQ 是斜槽，QR 为水平槽。实验时先使 A 球从斜槽上某一固定位置 G 由静止开始滚下，落到位于水平地面的记录纸上，留下痕迹，重复上述操作 10 次，得到 10 个落点痕迹。再把 B 球放在水平槽末端 R 在记录纸上的垂直投影点。B 球落点痕迹如图\ref{2000全国旧课程14b}所示，其中米尺水平放置，且平行于 G、R、O 所在的平面，米尺的零点 O 点对齐。
\twopicture[labelA=2000全国旧课程14a, labelB=2000全国旧课程14b, numA=甲, numB=乙, hA=4.0cm, hB=4.0cm]
{figs/14a.png}{figs/14b.png}
\begin{enumerate}
\item
碰撞后 B 球的水平射程应取为 \tkanswer[2cm]{64.7} $\Ucm$。

\item
在以下选项中，哪些是本次实验必须进行的测量？答：\tkanswer[2cm]{ABD}（填选项号）。

\fivechoices[answer=ABD]
{水平槽上未放 B 球时，测量 A 球落点位置到 O 点的距离}
{A 球与 B 球碰撞后，测量 A 球落点位置到 O 点的距离}
{测量 A 球或 B 球的直径}
{测量 A 球和 B 球的质量（或两球质量之比）}
{测量 G 点相对于水平槽面的高度}
\end{enumerate}

%% answer:
\jdanswer{
\begin{enumerate}
	\item 
	$64.7$
	
	\item
	ABD
\end{enumerate}
}

%% memo:
\memoanswer{
（1）图中画出了 B 球的 10 个落点，实验中应取平均位置。方法是用最小的圆将所有点圈在里面，圆心位置即为落点平均位置；找准平均位置，读数时应在刻度尺的最小刻度后面再估读一位。答案 64.7（3 分，答数在 64.2 到 65.2 范围内的都给分。）

（2）A、B 两球都从 O 点开始做平抛运动，且两球平抛时间相同，以平抛时间作为时间单位，则平抛的水平距离在数值上等于平抛初速度。设 A 未碰 B 时平抛水平位移 $s_{A}$，A、B 相碰后，A、B 两球的水平位移分别为 $s_{A}'$、$s_{B}'$，A、B 质量分别为 $m_{A}$、$m_{B}$，则碰前 A 的动量可写成 $m_{A}s_{A}$，碰后 A、B 总动量为 $m_{A}s_{A}' + m_{B}s_{B}'$。要验证动量是否守恒，即验证以上两动量是否相等，所以该实验应测量的物理量有 $m_{A}$、$m_{B}$、$s_{A}'$、$s_{B}'$，故 A、B、D 正确。答案 ABD（3 分，不是 ABD 的均给零分。）
}

\item
%% number: 15
%% paperName: 2000年全国高考旧课程第 15 题 6 分
%% typeId: 4
%% type: 实验题
%% chapter: 光学
%% point: 透镜
%% method: 无
%% score: 6
%% degree: 450
%% duplicateId: 0
%% body:
如图，一光源位于金属圆筒内部轴线上 A 点，与筒 B 端的距离为 $d$，$d$ 无法直接测量。另有凸透镜、光屏、米尺及带支架的光具座。现用这些器材测量 $d$。为此，先将圆筒、凸透镜、光屏依次放在光具座支架上，令圆筒轴线与透镜主光轴重合，屏与光源的距离足够远，使得移动透镜时，可在屏上两次出现光源的像，将圆筒及光屏位置固定。由光路的可逆性可知，第一次成像的像距等于第二次成像的物距。然后进行以下的测量：\tkanswer[10cm]{移动透镜，当屏上得到光源清晰的像时，测量像距 $v$；继续移动透镜，当屏上得到光源的另一个清晰的像时，测量端面 B 与透镜的距离 $l$}。

用测得的物理量可得 $d=$ \tkanswer[2cm]{$v-l$}。（应说明各符号所代表的物理量）
\onepicture[w=7cm, align=right]{figs/15.png}

%% answer:
\jdanswer{
\begin{enumerate}
	\item 
	移动透镜，当屏上得到光源清晰的像时，测量像距 $v$；继续移动透镜，当屏上得到光源的另一个清晰的像时，测量端面 B 与透镜的距离 $l$。
	
	\item
	$d=v-l$
\end{enumerate}
}

%% memo:
\memoanswer{
移动透镜，当屏上得到光源清晰的像时，测量像距 $v$；继续移动透镜，当屏上得到光源的另一个清晰的像时，测量端面 B 与透镜的距离 $l$。

$d = v - l$（6 分。完全答对给 6 分，否则不给分。如果测量方法与参考答案不同，但结果正确，同样给分。）
}

\item
%% number: 16
%% paperName: 2000年全国高考旧课程第 16 题 8 分
%% typeId: 4
%% type: 实验题
%% chapter: 电路
%% point: 电路实验
%% method: 无
%% score: 8
%% degree: 450
%% duplicateId: 0
%% body:
从下表中选出适当的实验器材，设计一电路来测量电流表 $A_{1}$ 的内阻 $r_{1}$，要求方法简捷，有尽可能高的测量精度，并能测得多组数据。

\begin{center}
\begin{tblr}{|c|l|}
\hline
器材（代号） & 规格 \\
\hline
电流表（$A_{1}$） & 量程 $10\UmA$，内阻 $r_{1}$ 待测（约 $40\UO$） \\
电流表（$A_{2}$） & 量程 $500\UuA$，内阻 $r_{2}=750\UO$ \\
电压表（V） & 量程 $10\UV$，内阻 $r_{3}=10\UkO$ \\
电阻（$R_{1}$） & 阻值约 $100\UO$，作保护电阻用 \\
滑动变阻器（$R_{2}$） & 总阻值约 $50\UO$ \\
电池（$E$） & 电动势 $1.5\UV$，内阻很小 \\
电键（K） & \\
导线若干 & \\
\hline
\end{tblr}
\end{center}

（1）在虚线方框中画出电路图，标明所用器材的代号。

（2）若选测量数据中的一组来计算，则所用的表达式为 $r_{1}=$ \tkanswer[3cm]{$\frac{I_{2}}{I_{1}}r_{2}$}，式中各符号的意义是：\tkanswer[8cm]{$I_{1}$ 表示通过电流表 $A_{1}$ 的电流，$I_{2}$ 表示通过电流表 $A_{2}$ 的电流，$r_{2}$ 表示电流表 $A_{2}$ 的内阻}。

%% answer:
\jdanswer{
\begin{enumerate}
	\item 
	如图所示
	
	\item
	$r_{1}=\frac{I_{2}}{I_{1}}r_{2}$
\end{enumerate}
}

%% memo:
\memoanswer{
（1）测量电路如图所示（5 分。若电阻 $R_{1}$ 与电池串联，不扣分；若 $R_{2}$ 用作限流电阻，正确的同样给分。）

\onepicture[w=4.2cm]{figs/16_an.png}

（2）直接判定必选的器材是电池 $E$、开关 $K$、导线若干。由 $A_{1}$ 量程是 $10\UmA$，$r_{1}$ 约为 $40\UO$，$E$ 电动势 $1.5\UV$，可判定必须再选电阻 $R_{1}$ 作为限流电阻，才保证通过 $A_{1}$ 的电流不至于超过 $10\UmA$。

要测 $A_{1}$ 的电阻，通过 $A_{1}$ 的电流值 $I_{1}$ 可由 $A_{1}$ 读出，只要能测出 $A_{1}$ 两端的电压 $U_{1}$，则可求出 $r_{1}=\frac{U_{1}}{I_{1}}$。

测电压要用电压表，但是若将电压表与 $A_{1}$ 并联，电流表 $A_{1}$ 上的电压最多是 $U_{1}=\frac{E}{R_{1}+r_{1}}\cdot r_{1}=\frac{1.5}{100+40}\times 40\UV = 0.43\UV$，而电压表的量程是 $10\UV$，对 $0.43\UV$ 的电压来说，电压表的指针只是稍微偏转了一点点，这样测得阻值误差太大，不符合题目要求，因而不能用电压表测电压。

由题意知电流表 $A_{2}$ 的内阻已知，通过电流表 $A_{2}$ 的电流 $I_{2}$ 可以算出电流表 $A_{2}$ 两端的电压，因而 $A_{2}$ 两端的电压 $U_{2}=I_{2}r_{2}$。

若将 $A_{1}$ 与 $A_{2}$ 并联，则 $U_{2}$ 即 $A_{1}$ 两端的电压，由此可以求出 $A_{1}$ 的内阻 $r_{1}=\frac{U_{2}}{I_{1}}=\frac{I_{2}r_{2}}{I_{1}}$，式中 $I_{1}$ 是 $A_{2}$、$A_{1}$ 并联时通过 $A_{1}$ 的电流，滑动变阻器 $R_{2}$ 限流、分压接法均可。

（2）$r_{1}=\frac{I_{2}}{I_{1}}r_{2}$，$I_{1}$ 表示通过电流表 $A_{1}$ 的电流，$I_{2}$ 表示通过电流表 $A_{2}$ 的电流，$r_{2}$ 表示电流表 $A_{2}$ 的内阻（3 分，未说明 $r_{2}$ 的不扣分。）
}

\item
%% number: 17
%% paperName: 2000年全国高考旧课程第 17 题 11 分
%% typeId: 6
%% type: 计算题
%% chapter: 气体性质
%% point: 气体连接体
%% method: 无
%% score: 11
%% degree: 450
%% duplicateId: 0
%% body:
一横截面积为 $S$ 的气缸水平放置，固定不动，气缸壁是导热的。两个活塞 A 和 B 将气缸分隔为 1、2 两气室，达到平衡时 1、2 两气室体积之比为 $3:2$，如图所示。在室温不变的条件下，缓慢推动活塞 A，使之向右移动一段距离 $d$。求活塞 B 向右移动的距离。不计活塞与气缸壁之间的摩擦。
\onepicture[w=4.5cm, align=right]{\includesvg{figs/17}}

%% answer:
\jdanswer{
$x=\frac{2}{5}d$
}

%% memo:
\memoanswer{
因气缸水平放置，又不计活塞的摩擦，故平衡时两气室内的压强必相等。设初态时气室内压强为 $p_{0}$，气室 1、2 的体积分别为 $V_{1}$ 和 $V_{2}$；在活塞 A 向右移动 $d$ 的过程中活塞 B 向右移动的距离为 $x$；最后气缸内压强为 $p$。因温度不变，分别对气室 1 和 2 的气体运用玻意耳定律，得

气室 1：$p_{0}V_{1}=p(V_{1}-Sd+Sx)$ ①

气室 2：$p_{0}V_{2}=p(V_{2}-Sx)$ ②

由①、②两式解得 $x=\frac{V_{2}}{V_{1}+V_{2}}d$ ③

由题意 $\frac{V_{1}}{V_{2}}=\frac{3}{2}$，得 $x=\frac{2}{5}d$ ④

评分标准：本题 11 分，①式 4 分，②式 3 分，得出③式 2 分，算出④式再给 2 分。
}

\item
%% number: 18
%% paperName: 2000年全国高考旧课程第 18 题 12 分
%% typeId: 6
%% type: 计算题
%% chapter: 交流电
%% point: 交流电
%% method: 无
%% score: 12
%% degree: 550
%% duplicateId: 0
%% body:
一小型发电机内的矩形线圈在匀强磁场中以恒定的角速度 $\omega$ 绕垂直于磁场方向的固定轴转动。线圈匝数 $n=100$。穿过每匝线圈的磁通量 $\Phi$ 随时间按正弦规律变化，如图所示。发电机内阻 $r=5.0\UO$，外电路电阻 $R=95\UO$。已知感应电动势的最大值 $E_{\mathrm{m}}=n\omega\Phi_{\mathrm{m}}$，其中 $\Phi_{\mathrm{m}}$ 为穿过每匝线圈磁通量的最大值。求串联在外电路中的交流电流表（内阻不计）的读数。
\onepicture[w=6.5cm, align=right]{figs/18.png}

%% answer:
\jdanswer{
$I=1.4\UA$
}

%% memo:
\memoanswer{
已知感应电动势的最大值

$E_{\mathrm{m}} = n\omega\Phi_{\mathrm{m}}$ ①

设线圈在磁场中转动的周期为 $T$，则有

$\omega = \frac{2\pi}{T}$ ②

根据欧姆定律，电路中电流的最大值为

$I_{\mathrm{m}} = \frac{E_{\mathrm{m}}}{R + r}$ ③

设交流电流表的读数为 $I$，它是电流的有效值，根据有效值与最大值的关系，有

$I = \frac{1}{\sqrt{2}} I_{\mathrm{m}}$ ④

由题给的 $\Phi-t$ 图线可读得

$\Phi_{\mathrm{m}} = 1.0 \times 10^{-2}\UWb$ ⑤

$T = 3.14 \times 10^{-2}\Us$ ⑥

解以上各式，并代入数据，得

$I = 1.4\UA$ ⑦

评分标准：本题 12 分，②式 1 分，③式 2 分，④式 3 分，⑤式 2 分，⑥式 2 分，⑦式 2 分。
}

\item
%% number: 19
%% paperName: 2000年全国高考旧课程第 19 题 13 分
%% typeId: 1
%% type: 单选题
%% chapter: 直线运动
%% point: 匀速直线运动
%% method: 分情况讨论
%% score: 13
%% degree: 450
%% duplicateId: 0
%% body:
如图所示，一辆实验小车可沿水平地面（图中纸面）上的长直轨道匀速向右运动。有一台发出细激光束的激光器装在小转台 M 上，到轨道的距离 MN 为 $d=10\Um$，转台匀速转动，使激光束在水平面内扫描，扫描一周的时间为 $T=60\Us$。光束转动方向如图中箭头所示。当光束与 MN 的夹角为 $45^\circ$ 时，光束正好射到小车上。如果再经过 $\Delta t=2.5\Us$，光束又射到小车上，则小车的速度为多少（结果保留两位有效数字）？
\onepicture[w=5.5cm, align=right]{figs/19.png}

%% answer:
\jdanswer{
$v_{1}=1.7\Ums$，$v_{2}=2.9\Ums$
}

%% memo:
\memoanswer{
在 $\Delta t$ 内，光束转过角度 $\Delta\phi=\frac{\Delta t}{T}\times360^{\circ}=15^{\circ}$ ①

由题意知，有两种可能，如图所示。

（1）光束照射小车时，小车正在接近 N 点，$\Delta t$ 内光束与 MN 的夹角从 $45^{\circ}$ 变为 $30^{\circ}$，小车走过 $L_{1}$，小车的速度为 $v_{1}=\frac{L_{1}}{\Delta t}$ ②，由几何关系知 $L_{1}=d(\tan45^{\circ}-\tan30^{\circ})$ ③，由②、③两式并代入数值，得 $v_{1}=1.7\Ums$ ④。

（2）光束照到小车时，小车正在远离 N 点，$\Delta t$ 内光束与 MN 的夹角从 $45^{\circ}$ 变为 $60^{\circ}$，小车走过 $L_{2}$，速度为 $v_{2}=\frac{L_{2}}{\Delta t}$ ⑤，由图可知 $L_{2}=d(\tan60^{\circ}-\tan45^{\circ})$ ⑥，由⑤、⑥两式并代入数值，得 $v_{2}=2.9\Ums$ ⑦。

\onepicture[w=7cm]{figs/19_an.png}

评分标准：本题 13 分。①式 2 分，②式 2 分，③式 3 分，④式 2 分，⑥式 2 分，⑦式 2 分，只考虑一种情形且正确的，只给 9 分。
}

\item
%% number: 20
%% paperName: 2000年全国高考旧课程第 20 题 12 分
%% typeId: 6
%% type: 计算题
%% chapter: 曲线运动
%% point: 人造地球卫星
%% method: 无
%% score: 12
%% degree: 350
%% duplicateId: 0
%% body:
2000 年 1 月 26 日我国发射了一颗同步卫星，其定点位置与东经 $98^\circ$ 的经线在同一平面内。若把甘肃省嘉峪关处的经度和纬度近似取为东经 $98^\circ$ 和北纬 $\alpha=40^\circ$，已知地球半径 $R$、地球自转周期 $T$、地球表面重力加速度 $g$（视为常量）和光速 $c$。试求该同步卫星发出的微波信号传到嘉峪关处的接收站所需的时间（要求用题给的已知量的符号表示）。

%% answer:
\jdanswer{
$t=\frac{\sqrt{\left(\frac{R^{2}gT^{2}}{4\pi^{2}}\right)^{2/3}+R^{2}-2R\left(\frac{R^{2}gT^{2}}{4\pi^{2}}\right)^{1/3}\cos\alpha}}{c}$
}

%% memo:
\memoanswer{
设 $m$ 为卫星质量，$M$ 为地球质量，$r$ 为卫星到地球中心的距离，$\omega$ 为卫星绕地心转动的角速度，由万有引力定律和牛顿定律有

$G\frac{mM}{r^{2}}=mr\omega^{2}$ ①

式中 $G$ 为万有引力恒量，因同步卫星绕地心转动的角速度 $\omega$ 与地球自转的角速度相等，有

$\omega=\frac{2\pi}{T}$ ②

因 $G\frac{Mm}{R^{2}}=mg$，得

$GM=gR^{2}$ ③

设嘉峪关到同步卫星的距离为 $L$，如图所示，由余弦定理

$L=\sqrt{r^{2}+R^{2}-2rR\cos\alpha}$ ④

所求时间为

$t=\frac{L}{c}$ ⑤

由以上各式得

$t=\frac{\sqrt{\left(\frac{R^{2}gT^{2}}{4\pi^{2}}\right)^{2/3}+R^{2}-2R\left(\frac{R^{2}gT^{2}}{4\pi^{2}}\right)^{1/3}\cos\alpha}}{c}$ ⑥

\onepicture[w=6cm]{figs/20_an.png}

评分标准：本题 12 分，①式 1 分，②式 2 分，③式 1 分，④式 5 分，⑤式 1 分，⑥式 2 分。
}

\item
%% number: 21
%% paperName: 2000年全国高考旧课程第 21 题 13 分
%% typeId: 6
%% type: 计算题
%% chapter: 磁场
%% point: 带电粒子在复合场中的运动
%% method: 无
%% score: 13
%% degree: 450
%% duplicateId: 0
%% body:
如图，两个共轴的圆筒形金属电极，外电极接地，其上均匀分布着平行于轴线的四条狭缝 a、b、c 和 d，外筒的外半径为 $r_{0}$。在圆筒之外的足够大区域中有平行于轴线方向的均匀磁场，磁感应强度的大小为 $B$。在两极间加上电压，使两圆筒之间的区域内有沿半径向外的电场。一质量为 $m$、带电量为 $+q$ 的粒子，从紧靠内筒且正对狭缝 a 的 S 点出发，初速度为零。如果该粒子经过一段时间的运动之后恰好又回到出发点 S，则两极之间的电压 $U$ 应是多少？（不计重力，整个装置在真空中。）
\onepicture[w=4.5cm, align=right]{figs/21.png}

%% answer:
\jdanswer{
$U=\frac{qr_{0}^{2}B^{2}}{2m}$
}

%% memo:
\memoanswer{
带电粒子从 S 出发，在两筒之间的电场力作用下加速，沿径向穿出 a 而进入磁场区，在洛伦兹力作用下做匀速圆周运动。粒子再回到 S 点的条件是能沿径向穿过狭缝 b。只要穿过了 b，粒子就会在电场力作用下先减速，再反向加速，经 b 重新进入磁场区。然后，粒子将以同样方式经过 c、d，再经过 a 回到 S 点。

设粒子射入磁场区的速度为 $v$，根据能量守恒，有

$\frac{1}{2}mv^{2}=qU$ ①

设粒子在洛伦兹力作用下做匀速圆周运动的半径为 $R$，由洛伦兹力公式和牛顿定律得

$m\frac{v^{2}}{R}=qBv$ ②

由前面分析可知，要回到 S 点，粒子从 a 到 b 必经过 $\frac{3}{4}$ 圆周，所以半径 $R$ 必定等于筒的外半径 $r_{0}$，即

$R = r_{0}$ ③

由以上各式解得

$U = \frac{qr_{0}^{2}B^{2}}{2m}$ ④

\onepicture[w=4.5cm]{figs/21_an.png}

评分标准：本题 13 分。①式 2 分，②式 2 分，经分析得出③式 6 分，解得④式 3 分。
}

\item
%% number: 22
%% paperName: 2000年全国高考旧课程第 22 题 14 分
%% typeId: 6
%% type: 计算题
%% chapter: 运动和力
%% point: 动量和能量
%% method: 无
%% score: 14
%% degree: 350
%% duplicateId: 0
%% body:
在原子核物理中，研究核子与核关联的最有效途径是“双电荷交换反应”。这类反应的前半部分过程和下述力学模型类似。两个小球 A 和 B 用轻质弹簧相连，在光滑的水平直轨道上处于静止状态。在它们左边有一垂直于轨道的固定挡板 P，右边有一小球 C 沿轨道以速度 $v_{0}$ 射向 B 球，如图所示。C 与 B 发生碰撞并立即结成一个整体 D。在它们继续向左运动的过程中，当弹簧长度变到最短时，长度突然被锁定，不再改变。然后，A 球与挡板 P 发生碰撞，碰后 A、D 都静止不动，A 与 P 接触而不粘连。过一段时间，突然解除锁定（锁定及解除锁定均无机械能损失）。已知 A、B、C 三球的质量均为 $m$。
\begin{enumerate}
\item
求弹簧长度刚被锁定后 A 球的速度。

\item
求在 A 球离开挡板 P 之后的运动过程中，弹簧的最大弹性势能。
\end{enumerate}
\onepicture[w=8cm, align=right]{figs/22.png}

%% answer:
\jdanswer{
\begin{enumerate}
	\item 
	$v_{2}=\frac{1}{3}v_{0}$
	
	\item
	$E_{p}'=\frac{1}{36}mv_{0}^{2}$
\end{enumerate}
}

%% memo:
\memoanswer{
（1）设 C 球与 B 球粘结成 D 时，D 的速度为 $v_{1}$，由动量守恒，有

$mv_{0} = (m + m)v_{1}$ ①

当弹簧压至最短时，D 与 A 的速度相等，设此速度为 $v_{2}$，由动量守恒，有

$2mv_{1} = 3mv_{2}$ ②

由①、②两式得 A 的速度

$v_{2}=\frac{1}{3}v_{0}$ ③

（2）设弹簧长度被锁定后，贮存在弹簧中的势能为 $E_{p}$，由能量守恒，有

$\frac{1}{2}\cdot2mv_{1}^{2}=\frac{1}{2}\cdot3mv_{2}^{2}+E_{p}$ ④

撞击 P 后，A 与 D 的动能都为零，解除锁定后，当弹簧刚恢复到自然长度时，势能全部转变成 D 的动能，设 D 的速度为 $v_{3}$，则有

$E_{p}=\frac{1}{2}(2m)\cdot v_{3}^{2}$ ⑤

当弹簧伸长，A 球离开挡板 P，并获得速度。当 A、D 的速度相等时，弹簧伸至最长。设此时的速度为 $v_{4}$，由动量守恒，有

$2mv_{3}=3mv_{4}$ ⑥

当弹簧伸到最长时，其势能最大，设此势能为 $E_{p}'$，由能量守恒，有

$\frac{1}{2}\cdot2mv_{3}^{2}=\frac{1}{2}\cdot3mv_{4}^{2}+E_{p}'$ ⑦

解以上各式得

$E_{p}'=\frac{1}{36}mv_{0}^{2}$ ⑧

评分标准：本题 14 分。第（1）问 5 分，其中①式 2 分，②式 2 分，③式 1 分；第（2）问 9 分，其中④式 2 分，⑤式 3 分，⑥式 1 分，⑦式 2 分，⑧式 1 分。
}

\end{enumerate}

\end{document}
